\section{Introduction}\label{sec:intro}

In an encrypted mempool, block builders order transactions before they can read them, which removes their ability to front-run or sandwich a transaction based on its contents. The hard question is who opens the transactions afterwards. There are three answers in current designs, and each has a cost.
\begin{itemize}[leftmargin=1.5em]
  \item \textbf{A committee} holds a threshold decryption key and opens what was ordered, as in Shutter~\cite{Shutter}. Privacy then rests on no threshold of its members colluding, and liveness on enough of them being online.
  \item \textbf{A trusted execution environment} opens transactions inside hardware whose vendor and side channels must be trusted.
  \item \textbf{The sender} opens its own transaction after the ordering, as in sealed transactions and the self-decryption path of LUCID~\cite{LUCID}. No one else is trusted, but the sender sees the ordering first and can withhold its key when the ordering does not suit it: a free option.
\end{itemize}
Hybrid designs~\cite{WM25} let either the sender or a committee open, which removes the free option while the committee is online and the committee's liveness problem while the sender is honest. Their bad case remains a dishonest sender while the committee is offline.

\leanseal makes the free option observable and penalizable without any committee. The sender opens its own transaction in the normal case, but it attaches a time-lock puzzle that anyone can solve to recover the key after a fixed amount of sequential computation, about 30 seconds on Ethereum. For a well-formed puzzle, withholding the key therefore cannot keep a transaction hidden, though it can still delay or cancel its execution; a deposit forfeited on a missed reveal must deter that. This is a timed commitment~\cite{BN00}. The puzzle uses only SHA-256: it needs no trusted setup, no committee and no hardware trust, and the idealized version is secure against quantum openers. The price is the sender's work. The puzzle is the chained-hash puzzle of Mahmoody, Moran and Vadhan~\cite{MMV11}, in which the sender computes as many hashes as an opener, only in parallel, and they proved that no hash-only puzzle can do better. This paper asks whether that price is affordable for sealing transactions, and what an encrypted mempool built on it can and cannot provide.

\begin{table}[t]
  \centering\small
  \caption{The reference puzzle at a glance: a 30-second lock of $n = 16$ chains of $2^{29}$ SHA-256 steps with 24-byte links. Times are on an Apple M4 MacBook Air (\cref{sec:costs}); building and forced opening are half our measurements of the one-minute puzzle ($m = 2^{30}$, run with 32-byte links and the same chain work), since both are linear in the delay.}
  \label{tab:glance}
  \begin{tabularx}{\textwidth}{@{}Xrl@{}}
    \toprule
    & Value & Basis \\
    \midrule
    Earliest opening by an attacker & 30\,s & \tagm{assumed}: 3.5\,ns per chained hash \\
    Opening when the sender reveals the key, the normal case & 12.0\,ms & \tagm{measured}: command line, end to end \\
    Forced opening without the sender, one laptop core & 3\,min 21\,s & \tagm{measured}, halved \\
    Building the puzzle, before the transaction exists & 21\,s & \tagm{measured}, halved: all CPU cores \\
    Added to a transaction, on-demand or stockpiled puzzle & 512 or 152\,B & exact \\
    \bottomrule
  \end{tabularx}
\end{table}

\paragraph{Our results.}
\begin{enumerate}[label=\arabic*.,leftmargin=1.5em]
  \item \textbf{A design.} We describe \leanseal as a decryption path for an encrypted mempool: self-decryption with a time-lock fallback (\cref{sec:mempool}). It uses the same data as LUCID and the hybrid design, a ciphertext bound to a commitment to its key, so its puzzle fits them as a fallback that needs no third party.
  \item \textbf{A concrete security proof.} In the random-oracle model, a classical attacker derives the key even one oracle round early with probability at most about $2^{-117}$ for a 30-second lock against the hash rate of the whole Bitcoin network (\cref{sec:security}). Turning rounds into seconds needs one hardware number, the fastest time for one chained SHA-256, which we price at 3.5\,ns, 3.7 times faster than the fastest published design we found.
  \item \textbf{Affordable costs.} A 30-second puzzle, enough when the order is fixed within a slot, takes about 21\,s to build on a laptop CPU and can be built before the transaction exists (\cref{sec:choosing-t} chooses the delay per use case). A sender that keeps a stock can use single-chain puzzles, which add 152 bytes to a transaction (\cref{sec:costs}).
  \item \textbf{An honest account of the trade-offs.} A time lock also reveals transactions that are never included; we argue that this holds for every design whose fallback opening is input-free, and show how validity windows, private submission and a committee limit it. Withholding can still delay or cancel execution, so a deposit carries the economics; the reveal deadline makes a missed reveal observable, and the time lock separately ensures that a well-formed puzzle is eventually disclosed (\cref{sec:tradeoffs}). We compare \leanseal with threshold, TEE, hybrid and delay-encryption designs (\cref{sec:comparison}).
  \item \textbf{An implementation}: a specified format with a Rust implementation, a second implementation in Python, test vectors and benchmarks, and a prototype on an EVM chain that demonstrates sealed submission, key revelation, ordered execution within a batch and forced disclosure; its searcher sandwiches only plaintext swaps, for a 9.25\% loss, so it does not yet demonstrate sandwich protection (\cref{sec:poc})~\cite{leanSealRepo}.
\end{enumerate}
Neither the puzzle nor its implementation is new in itself (\cref{sec:related}). What is new is a concrete bound for an exact encoding, measured costs against an explicit attacker, and its analysis as an encrypted mempool.
